\documentclass{article}
\usepackage[utf8]{inputenc}
\usepackage{amsmath}
\usepackage{listings}
\usepackage{xcolor}
\usepackage{hyperref}
\usepackage{geometry}
\geometry{a4paper, margin=1in}

\title{Detailed Report: Code Analysis of \texttt{1\_task.R}}
\author{}
\date{\today}

\definecolor{codegreen}{rgb}{0,0.6,0}
\definecolor{codegray}{rgb}{0.5,0.5,0.5}
\definecolor{codepurple}{rgb}{0.58,0,0.82}
\definecolor{backcolour}{rgb}{0.95,0.95,0.92}

\lstdefinestyle{mystyle}{
    backgroundcolor=\color{backcolour},   
    commentstyle=\color{codegreen},
    keywordstyle=\color{magenta},
    numberstyle=\tiny\color{codegray},
    stringstyle=\color{codepurple},
    basicstyle=\ttfamily\footnotesize,
    breakatwhitespace=false,         
    breaklines=true,                 
    captionpos=b,                    
    keepspaces=true,                 
    numbers=left,                    
    numbersep=5pt,                  
    showspaces=false,                
    showstringspaces=false,
    showtabs=false,                  
    tabsize=2
}
\lstset{style=mystyle}

\begin{document}

\maketitle

\section{Introduction}
This report provides a detailed, line-by-line analysis of the \texttt{1\_task.R} file, which corresponds to Task 1 of the project. The primary objective of the script is to estimate different Beta mixture models (ranging from 1 to 5 components) using Bayesian methods via JAGS and to select the optimal model based on the WAIC (Watanabe-Akaike Information Criterion).

\section{Line-by-Line Code Analysis}

\subsection{Data and Structure Initialization}
\begin{lstlisting}[language=R]
# ==============================================================================
# TASK 1: Beta Mixture Models and WAIC Selection
# ==============================================================================

# JAGS requires data in a list format
jags_data <- list(
  y = y,
  N = N
)

# Lists to store the results
waic_list <- list()
models_mcmc <- list()
\end{lstlisting}
\begin{itemize}
  \item \textbf{Lines 1-5}: Introductory comments.
  \item \textbf{Lines 6-9}: A list named \texttt{jags\_data} is created, containing the observations (\texttt{y}) and the total number of observations (\texttt{N}). JAGS requires input data to be structured as a list. We assume that \texttt{y} and \texttt{N} were defined previously (e.g., in \texttt{0\_setup.R}).
  \item \textbf{Lines 11-14}: Two empty lists, \texttt{waic\_list} and \texttt{models\_mcmc}, are initialized. These are used to store the computed WAIC values and the MCMC (Markov Chain Monte Carlo) samples extracted for each estimated model, respectively.
\end{itemize}

\subsection{JAGS Model Fitting Function Definition}
\begin{lstlisting}[language=R, firstnumber=15]
# Function to run JAGS model for a given H
fit_beta_mixture <- function(H, data, n_iter = 500, n_burnin = 500) {
\end{lstlisting}
\begin{itemize}
  \item \textbf{Line 16}: The \texttt{fit\_beta\_mixture} function is defined. This function takes as input the number of mixture components \texttt{H}, the data \texttt{data}, and the MCMC parameters (\texttt{n\_iter} for sampling iterations and \texttt{n\_burnin} for the warm-up phase).
\end{itemize}

\subsubsection{\texorpdfstring{Model for $H = 1$ (Single Beta)}{}}
\begin{lstlisting}[language=R, firstnumber=18]
  if (H == 1) {
    # Single Beta model
    model_string <- "
    model {
      for (i in 1:N) {
        y[i] ~ dbeta(alpha, beta)
      }
      
      # Priors (weakly informative)
      alpha ~ dgamma(0.1, 0.1)
      beta ~ dgamma(0.1, 0.1)
    }
    "
    params_to_monitor <- c("alpha", "beta")
\end{lstlisting}
\begin{itemize}
  \item \textbf{Lines 18-32}: Handles the case where the number of components is 1, representing a standard, non-mixture Beta model.
  \item \textbf{Lines 22-24}: The model likelihood: each observation \texttt{y[i]} follows a Beta distribution with parameters \texttt{alpha} and \texttt{beta}.
  \item \textbf{Lines 27-28}: Prior distributions are defined for \texttt{alpha} and \texttt{beta}. Gamma(0.1, 0.1) distributions are used, which are weakly informative, allowing the data to drive the estimation.
  \item \textbf{Line 31}: Specifies the parameters we want JAGS to return and monitor: \texttt{alpha} and \texttt{beta}.
\end{itemize}

\subsubsection{\texorpdfstring{Model for $H > 1$ (Beta Mixture)}{}}
\begin{lstlisting}[language=R, firstnumber=33]
  } else {
    # Mixture of H Betas with ordering constraint on means to prevent label switching
    model_string <- paste0("
    model {
      for (i in 1:N) {
        y[i] ~ dbeta(alpha[z[i]], beta[z[i]])
        z[i] ~ dcat(pi[1:H])
      }
      
      # Priors for mixing weights
      pi[1:H] ~ ddirich(rep(1, H))
      
      # Priors for Beta parameters using mu and precision (kappa) 
      # with ordering on mu to prevent label switching
      for (h in 1:H) {
        mu_raw[h] ~ dunif(0, 1)
        kappa[h] ~ dgamma(0.1, 0.1)
      }
      
      # Sort mu_raw to ensure mu[1] < mu[2] < ... < mu[H]
      mu[1:H] <- sort(mu_raw)
      
      for (h in 1:H) {
        alpha[h] <- mu[h] * kappa[h]
        beta[h] <- (1 - mu[h]) * kappa[h]
      }
    }
    ")
    
    # Update data list to include H
    data$H <- H
    params_to_monitor <- c("alpha", "beta", "pi")
  }
\end{lstlisting}
\begin{itemize}
  \item \textbf{Lines 33-65}: If $H > 1$, the mixture model is defined.
  \item \textbf{Line 38}: Each observation $y_i$ follows a Beta distribution whose parameters depend on the assigned component $z_i$.
  \item \textbf{Line 39}: The latent variable $z_i$ (cluster assignment) follows a categorical distribution with probabilities given by the mixture weights \texttt{pi}.
  \item \textbf{Line 43}: A symmetric Dirichlet prior is used for the mixture weights \texttt{pi}, ensuring the weights sum to 1.
  \item \textbf{Lines 47-50}: To address the component identifiability issue (label switching), the Beta distributions are parameterized using the mean \texttt{mu} and precision \texttt{kappa}. Initially, raw means are generated from a Uniform(0,1) distribution, and precisions from a weakly informative Gamma.
  \item \textbf{Line 53}: The means are sorted in ascending order. This step is crucial for preventing label switching: it imposes a strict ordering on the components, ensuring the algorithm does not swap cluster labels during sampling.
  \item \textbf{Lines 55-58}: The canonical shape parameters of the Beta distribution, \texttt{alpha} and \texttt{beta}, are calculated from the sorted means \texttt{mu} and precisions \texttt{kappa}.
  \item \textbf{Lines 63-64}: $H$ is added to the data list, and JAGS is instructed to monitor the parameters \texttt{alpha}, \texttt{beta}, and the weights \texttt{pi}.
\end{itemize}

\subsubsection{JAGS Execution}
\begin{lstlisting}[language=R, firstnumber=67]
  # Write model to temporary file
  model_file <- tempfile()
  writeLines(model_string, con = model_file)
  
  # Initialize and run JAGS
  cat(paste("\nFitting model with H =", H, "...\n"))
  jags_model <- jags.model(file = model_file, data = data, n.chains = 2, n.adapt = n_burnin, quiet = TRUE)
  
  # Burn-in
  update(jags_model, n_burnin)
  
  # Sample
  samples <- coda.samples(jags_model, variable.names = params_to_monitor, n.iter = n_iter)
  
  return(samples)
}
\end{lstlisting}
\begin{itemize}
  \item \textbf{Lines 68-69}: The model string is written to a temporary file for JAGS to read.
  \item \textbf{Line 73}: The JAGS model is initialized by compiling the code, using 2 chains (\texttt{n.chains = 2}) and starting an initial adaptation phase of length \texttt{n\_burnin}.
  \item \textbf{Line 76}: A further burn-in of \texttt{n\_burnin} iterations is executed to allow the chains to converge toward the stationary distribution before collecting final samples.
  \item \textbf{Line 79}: Posterior samples are drawn for \texttt{n\_iter} iterations for the monitored variables. The result is returned by the function.
\end{itemize}

\subsection{Model Estimation and Evaluation (Main Loop)}
\begin{lstlisting}[language=R, firstnumber=84]
# Run models for H = 1, 2, 3, 4, 5
for (H in 1:5) {
  n_it <- if(exists("GLOBAL_N_ITER")) GLOBAL_N_ITER else 500
  n_bu <- if(exists("GLOBAL_N_BURNIN")) GLOBAL_N_BURNIN else 500
  
  samples <- fit_beta_mixture(H, jags_data, n_iter = n_it, n_burnin = n_bu)
  models_mcmc[[H]] <- samples
\end{lstlisting}
\begin{itemize}
  \item \textbf{Line 85}: A loop begins to test models with a number of components $H$ ranging from 1 to 5.
  \item \textbf{Lines 86-87}: The number of iterations and burn-in values are set. If global variables for these values exist (presumably defined in \texttt{0\_setup.R}), they are used; otherwise, it defaults to 500.
  \item \textbf{Lines 89-90}: The \texttt{fit\_beta\_mixture} function is called, and the returned samples are saved in the \texttt{models\_mcmc} list.
\end{itemize}

\subsubsection{Log-Likelihood Calculation for WAIC}
\begin{lstlisting}[language=R, firstnumber=92]
  # Extract and compute log_lik for WAIC
  samples_mat <- as.matrix(samples)
  S <- nrow(samples_mat)
  log_lik_mat <- matrix(0, nrow = S, ncol = N)
  
  cat("Computing WAIC...\n")
  for (s in 1:S) {
    if (H == 1) {
      log_lik_mat[s, ] <- dbeta(y, samples_mat[s, "alpha"], samples_mat[s, "beta"], log = TRUE)
    } else {
      lik_i <- rep(0, N)
      for (h in 1:H) {
        pi_h <- samples_mat[s, paste0("pi[", h, "]")]
        alpha_h <- samples_mat[s, paste0("alpha[", h, "]")]
        beta_h <- samples_mat[s, paste0("beta[", h, "]")]
        lik_i <- lik_i + pi_h * dbeta(y, alpha_h, beta_h, log = FALSE)
      }
      log_lik_mat[s, ] <- log(lik_i)
    }
  }
\end{lstlisting}
\begin{itemize}
  \item \textbf{Lines 93-95}: The MCMC object is converted into a matrix of $S$ rows (total samples across all chains). A \texttt{log\_lik\_mat} matrix of dimensions $S \times N$ (samples by observations) is initialized to hold the pointwise log-likelihood.
  \item \textbf{Line 98}: Iterates over each MCMC sample $s$.
  \item \textbf{Line 100}: If $H=1$, the log-likelihood for each observation is calculated in a vectorized manner using the $\alpha$ and $\beta$ parameters extracted from that sample.
  \item \textbf{Lines 101-110}: If $H > 1$, the observation likelihood is the weighted sum of the individual mixture component likelihoods. For each component $h$, it calculates $\pi_h \times \text{Beta}(y \mid \alpha_h, \beta_h)$ and sums them up. Finally, at line 109, the logarithm is applied to the total likelihood of the mixture.
\end{itemize}

\subsubsection{WAIC Calculation and Model Selection}
\begin{lstlisting}[language=R, firstnumber=113]
  # Calculate WAIC using loo package
  waic_res <- waic(log_lik_mat)
  waic_list[[H]] <- waic_res
  
  cat(sprintf("WAIC for H = %d: %.2f\n", H, waic_res$estimates["waic", "Estimate"]))
}

# Select best model based on WAIC
waic_values <- sapply(waic_list, function(w) w$estimates["waic", "Estimate"])
best_H <- which.min(waic_values)

cat(sprintf("\nBest model selected by WAIC has H = %d components (WAIC = %.2f).\n", best_H, waic_values[best_H]))
\end{lstlisting}
\begin{itemize}
  \item \textbf{Lines 114-115}: The \texttt{waic()} function (presumably from the \texttt{loo} package) takes the log-likelihood matrix \texttt{log\_lik\_mat} as input and calculates the WAIC index, then saves it in the \texttt{waic\_list}.
  \item \textbf{Lines 120-122}: After the loop concludes, the estimated WAIC value for each model is extracted, and the index (\texttt{best\_H}) corresponding to the model that minimizes the WAIC is found (a lower value indicates a better trade-off between goodness of fit and model complexity).
\end{itemize}

\subsection{Saving the Results}
\begin{lstlisting}[language=R, firstnumber=126]
# Save results for subsequent tasks
best_model <- best_H
waic_results <- waic_list
mcmc_samples <- models_mcmc[[best_H]]
\end{lstlisting}
\begin{itemize}
  \item \textbf{Lines 126-129}: The number of components for the optimal model, the complete WAIC results, and the MCMC samples of the selected best model are saved into specifically named variables (\texttt{best\_model}, \texttt{waic\_results}, \texttt{mcmc\_samples}). This ensures that this essential information is available in the global environment to be reused by subsequent scripts (e.g., \texttt{2\_task.R}).
\end{itemize}

\section{Conclusion}
The code is well-structured. A highly important feature is the ordering of the means (\texttt{mu}) in the JAGS block to prevent the "label switching" problem, an excellent practice in Bayesian modeling of mixtures. The pointwise log-likelihood calculation (necessary for WAIC) is correctly implemented by converting the JAGS samples into a matrix and extracting the log-likelihood pointwise for each observation and each MCMC iteration.

\end{document}
